% Current manuscript Table 1. Source: input_markdown/sections/02_background_and_scope.md.
% Evidence types describe the evidence examined; they are not capability scores.
\begin{table*}[t]
  \caption{Closely related research and the evidence it examines. The comparison identifies the study focus rather than claiming absent capabilities or empirical superiority.}
  \label{tab:related}
  \centering
  \small
  \setlength{\tabcolsep}{4pt}
  \renewcommand{\arraystretch}{1.12}
  \begin{tabularx}{\textwidth}{@{}>{\raggedright\arraybackslash}p{0.19\textwidth} >{\centering\arraybackslash}p{0.048\textwidth} >{\raggedright\arraybackslash}p{0.325\textwidth} >{\raggedright\arraybackslash}X@{}}
    \toprule
    Work & Type & Research focus & Focus of the present study \\
    \midrule
    FP-Scanner~\cite{ref04} & \HGTag{D} & Countermeasure-induced inconsistencies & Associated Native/Host references \\
    Picasso~\cite{ref05} & \HGTag{C} & Claimed device-class checking & Current-record semantic compatibility \\
    FP-Inconsistent~\cite{ref06} & \HGTag{I} & Evasive bot traffic & Controlled edits vs.\ normal changes \\
    BrowserFM~\cite{ref07} & \HGTag{G} & Configuration modeling and analysis & Applicability + evidence-availability costs \\
    WebView studies~\cite{ref08,ref09,ref10} & \HGTag{S} & Exposure + cross-boundary tracking & Report-integrity assessment \\
    \midrule
    \textbf{HybridGuard} & \HGTag{R} & Contextual references; explicit domains & Detection, normal alarms, unknown outcomes \\
    \bottomrule
  \end{tabularx}
  \par\vspace{2pt}
  \footnotesize
  \begin{tabularx}{\textwidth}{@{}>{\raggedright\arraybackslash}X >{\raggedright\arraybackslash}X@{}}
    \HGTag{D} Browser attributes; JavaScript/HTTP discrepancies & \HGTag{C} Canvas challenge responses \\
    \HGTag{I} Within-request and temporal inconsistencies & \HGTag{G} Fingerprint/configuration dependencies \\
    \HGTag{S} App, JavaScript, and runtime observations & \HGTag{R} Native, Host, App Web; Browser optional
  \end{tabularx}
\end{table*}
